% !TeX spellcheck = en_US
\addsection{Scenarios}{\skills/intelligence.png}

\begin{multicols*}{2}

Heroes III is a scenario-driven game. In the first step of setup of each game you choose a \textbf{Scenario} from the Mission Book. Scenarios have a backstory and determine the map layout, the starting conditions for players, and the requirements for winning the Scenario. Furthermore, Scenarios can also introduce entirely unique game rules.

\subheader{Types of Scenarios}

There are the following five basic types of Scenarios:

\subsection*{Clash}
A fully competitive mode.
Your opponents control Neutral Units in your Combat.

\subsection*{Campaign}
A single player mode\index{Solo Mode} of interconnected Scenarios against an enemy AI.
For rules unique to Solo Mode, see \pagelink{Campaign}{the Campaign rules}.
Further rules changes are detailed in the Campaign Mission Books.

\subsection*{Alliance}
Usually a 2 vs.~2 team-based mode with expanded opportunities to \pagelink{Trading}{trade} between allied players.

\subsection*{Cooperative}
A Cooperative mode where everyone shares the same goal. \pagelink{AI Units}{The AI rules} and expanded \pagelink{Trading}{Trading rules} apply here.

\subsection*{Tournament}
A 1 vs.~1 mode with a different way of building Map and scoring. See \pagelink{Tournament Mode}{Tournament rules}.

\bigskip
\begin{expansion}[before=\vspace*{-10mm}]{battlefield}
  \pagelink{Adventure and Skirmish Modes}{Battlefield Expansion} introduces two focused 1 vs.~1 game modes \textbf{which work without Scenarios}.

  \subsubsection*{Adventure}
  The Map exploration is replaced with Adventure cards, culminating in a tactical combat on a large hexagonal battlefield, see \pagelink{Adventure}{Adventure Mode} for details.

  \subsubsection*{Skirmish}
  Skip the exploration phase and engage directly in tactical combat, see \pagelink{Skirmish}{Skirmish Mode} for details.
\end{expansion}

\subheader{Scenario Difficulty}

    You can modify the Difficulty of a Scenario.
    There are four different Difficulties, each with a different starting bonus\index{Starting Bonus} that players receive during step 16 of the setup:
    \begin{itemize}
        \item \textbf{Easy} – Roll 2 Resource Dice \svg{resource_die} and receive Resources from both – OR – \textbf{Search} (2) the Artifact Deck, twice.
        \item \textbf{Normal} – Roll 2 Resource Dice \svg{resource_die} and receive the Resources from one of them – OR – \textbf{Search} (2) the Artifact Deck.
        \item \textbf{Hard} – Roll 1 Resource Die \svg{resource_die} and receive the Resources on it – OR – reveal cards from the top of the Artifact Deck until you find 1 Minor Artifact and add it to your hand.
        \item \textbf{Impossible} – No starting bonus.
    \end{itemize}
    \columnbreak

    Campaign Scenarios have unique bonuses that \textbf{replace} the regular starting bonus.

    \pagelink{deck-building}{Searching \textbf{Artifacts}} is explained in general card rules. Artifacts received from a starting bonus should be placed into your \textbf{hand} and not shuffled into your Starting Deck.
    If you searched for any Artifacts, shuffle the Artifact Deck and its Discard Pile\index{Discard Pile} together afterwards, and then discard one Artifact from the top to form the Artifact Discard Pile again.\par
    The chosen difficulty also determines the number and type of neutral enemies that are encountered during Neutral Combat according to the \hyperlink{Difficulty Table}{table} at the back cover of the book.

    \subheader{Victory conditions}
    All Scenarios have their victory conditions\index{Victory Conditions} described in the Scenario Book. The most common victory conditions require defeating certain enemy Towns, Heroes, or Neutral Units; conquering the Holy Grail; or controlling certain or most Locations over time.

    In addition, it is always possible to be \textbf{Eliminated} from any Scenario in the following ways:
    \begin{itemize}
        \item Play 3 full Rounds without controlling a Town or a Settlement.
        Count the number of Rounds left using any suitable component.
        \item Lose Combat with your \textbf{Main Hero} when you have no Towns or Settlements left, including when defending your last Town or Settlement.
    \end{itemize}
    Eliminated players are immediately removed from the game.
    Discard their Faction Cubes and Hero models from the Map.
    Treat the cards in their Deck as being removed from the game for the rest of the Scenario.
    If you are Eliminated, you may still participate in the game by controlling Neutral Units.\par
    \note{4}{If you Eliminate all enemy Factions, you immediately win the Scenario.}\par

\columnbreak

    In \textbf{Clash} Scenarios with three or more players, collecting \textbf{a Faction Cube} from every enemy player immediately wins you the game.
    You get such Faction Cubes by defeating Main Heroes or conquering enemy Towns.
    Other Scenario specific rules may also modify the outcome of collecting Faction Cubes.

    After finishing a \textbf{Solo Campaign} Scenario, see the \pagelink{Campaign}{Campaign section} how to prepare for the next Scenario of the Campaign.

    \vspace*{\fill}
    \hspace{-3em}
    {\transparent{0.7}{\includegraphics[width=\iftoggle{printable}{1.5}{1.7}\linewidth]{\art/medusa.png}}}

\end{multicols*}